\begin{helper}
Fix \(n,m,k\), an edge parameter \(p\), parameters
\(\varepsilon,\varepsilon_1,\varepsilon_2\), and a two-bite sample
\(\omega\).  Put
\[
A_0=k\,n^{1/4-\varepsilon},\qquad
L=3(\log n)^2.
\]
Assume \(\noderef{FiberAndDegreeEvent}(\omega,\varepsilon_2)\), assume that
\(\noderef{TwoBiteMediumClosedPairsEvent}(k,\varepsilon,\varepsilon_1,\omega)\)
fails, and assume the deterministic numerical inputs
\[
A_0\ge0,\qquad \noderef{TwoBiteLargeCutoff}(\varepsilon,n)>0,\qquad
\noderef{TwoBiteLargeCutoff}(\varepsilon,n)A_0
  \le \varepsilon_1 k^2,
\]
\[
(1+\varepsilon_2)(\log n)^2\le 3(\log n)^2,\qquad L>0 .
\]
Then there is a \(k\)-set \(I\subseteq \mathrm{Fin}(n)\) for which the medium
closed-pairs bound fails, and there is a nonempty family
\(B\subseteq \noderef{TwoBiteMediumClass}(\omega,\varepsilon,I)\) such that
\[
A_0<\sum_{x\in B}|\noderef{TwoBiteX}(\omega,I,x)|
\le A_0+\noderef{TwoBiteLargeCutoff}(\varepsilon,n),
\]
\[
|B|\,\noderef{TwoBiteSmallCutoff}(\varepsilon,n)
\le A_0+\noderef{TwoBiteLargeCutoff}(\varepsilon,n).
\]
Write
\[
B_R=\{r:\operatorname{inl}(r)\in B\},\qquad
B_B=\{b:\operatorname{inr}(b)\in B\},
\]
\[
P_R=\noderef{TwoBiteRedProjection}(\omega)(I),\qquad
P_B=\noderef{TwoBiteBlueProjection}(\omega)(I).
\]
Let \(E^\rightarrow_{I,B}\) be the number of ordered present red candidate
pairs
\[
(r,s)\in B_R\times P_R
\quad\text{with}\quad
\noderef{TwoBiteRedGraph}(\omega).\operatorname{Adj}(r,s),
\]
plus the number of ordered present blue candidate pairs
\[
(b,c)\in B_B\times P_B
\quad\text{with}\quad
\noderef{TwoBiteBlueGraph}(\omega).\operatorname{Adj}(b,c).
\]
Let \(E^{\mathrm{coord}}_{I,B}\) be the corresponding number of actual
base-graph edge coordinates: in each color, normalize every present ordered
candidate \((x,y)\) to the nonloop coordinate
\((\min(x,y),\max(x,y))\), and count distinct normalized red coordinates plus
distinct normalized blue coordinates.  Then
\[
E^{\mathrm{coord}}_{I,B}>A_0/(2L)=A_0/(6(\log n)^2).
\]
\end{helper}
\begin{proof}
Unfold \(\noderef{TwoBiteMediumClosedPairsEvent}\).  Since the event fails,
there is a finite set \(I\subseteq\mathrm{Fin}(n)\) with \(|I|=k\) such that
\[
\neg\noderef{ClosedPairsBound}\left(
|\noderef{TwoBiteClosedPairsInClass}(\omega,I,
  \noderef{TwoBiteMediumClass}(\omega,\varepsilon,I))|,
  \varepsilon_1,k\right).
\]
Because \(\noderef{ClosedPairsBound}(z,\varepsilon_1,k)\) is the inequality
\(z\le \varepsilon_1 k^2\), this negation is the strict inequality
\[
\varepsilon_1k^2<
|\noderef{TwoBiteClosedPairsInClass}(\omega,I,
  \noderef{TwoBiteMediumClass}(\omega,\varepsilon,I))|. \tag{1}
\]

Let
\[
M_I=\{x:\noderef{TwoBiteMediumClass}(\omega,\varepsilon,I,x)\}.
\]
For every \(x\in M_I\), unfolding \(\noderef{TwoBiteMediumClass}\) gives
\[
\noderef{TwoBiteSmallCutoff}(\varepsilon,n)
<|\noderef{TwoBiteX}(\omega,I,x)|
\le \noderef{TwoBiteLargeCutoff}(\varepsilon,n). \tag{2}
\]
The definition \(\noderef{TwoBiteClosedPairsInClass}\) is the union, over
\(x\in M_I\), of the final pairs generated by
\(\noderef{TwoBiteX}(\omega,I,x)\).  Finite subadditivity of cardinality
therefore gives
\[
|\noderef{TwoBiteClosedPairsInClass}(\omega,I,M_I)|
\le \sum_{x\in M_I}\binom{|\noderef{TwoBiteX}(\omega,I,x)|}{2}. \tag{3}
\]
For each \(x\in M_I\), the upper side of (2) implies
\[
\binom{|\noderef{TwoBiteX}(\omega,I,x)|}{2}
\le
\noderef{TwoBiteLargeCutoff}(\varepsilon,n)
  |\noderef{TwoBiteX}(\omega,I,x)|. \tag{4}
\]
Indeed, the binomial coefficient is at most
\(|\noderef{TwoBiteX}(\omega,I,x)|^2/2\), and the nonnegative integer
\(|\noderef{TwoBiteX}(\omega,I,x)|\) is at most the large cutoff by (2).
Combining (1), (3), and (4) gives
\[
\varepsilon_1k^2<
\noderef{TwoBiteLargeCutoff}(\varepsilon,n)
\sum_{x\in M_I}|\noderef{TwoBiteX}(\omega,I,x)|.
\]
Using
\(\noderef{TwoBiteLargeCutoff}(\varepsilon,n)A_0\le\varepsilon_1k^2\) and
\(\noderef{TwoBiteLargeCutoff}(\varepsilon,n)>0\), division by the positive
large cutoff yields
\[
A_0<\sum_{x\in M_I}|\noderef{TwoBiteX}(\omega,I,x)|. \tag{5}
\]

For a subfamily \(C\subseteq M_I\), define
\[
w(x)=|\noderef{TwoBiteX}(\omega,I,x)|.
\]
Let \(C_R=\{r:\operatorname{inl}(r)\in C\}\) and
\(C_B=\{b:\operatorname{inr}(b)\in C\}\).  Let \(E^\rightarrow(C)\) be the
number of ordered present red pairs in \(C_R\times P_R\), plus the number of
ordered present blue pairs in \(C_B\times P_B\), with the same adjacency
conditions as in the statement.  We prove
\[
\sum_{x\in C}w(x)\le L\,E^\rightarrow(C). \tag{6}
\]

First consider a red incidence.  Thus \(x=\operatorname{inl}(r)\in C\) and
\(u\in\noderef{TwoBiteX}(\omega,I,x)\).  By the definition
\(\noderef{TwoBiteX}\), the vertex \(u\) lies in \(I\) and in
\(\noderef{TwoBiteLiftedNeighborhood}(\omega,\operatorname{inl}(r))\).
Unfolding \(\noderef{TwoBiteLiftedNeighborhood}\), this means
\[
\noderef{TwoBiteRedGraph}(\omega).\operatorname{Adj}
  \bigl(r,\noderef{TwoBiteRedProjection}(\omega,u)\bigr).
\]
Since \(u\in I\), the projection
\(\noderef{TwoBiteRedProjection}(\omega,u)\) lies in \(P_R\).  Hence the
ordered pair
\[
\bigl(r,\noderef{TwoBiteRedProjection}(\omega,u)\bigr)\in C_R\times P_R
\]
is one of the ordered red pairs counted by \(E^\rightarrow(C)\).  The blue
incidence case is identical: if \(x=\operatorname{inr}(b)\in C\) and
\(u\in\noderef{TwoBiteX}(\omega,I,x)\), then
\[
\bigl(b,\noderef{TwoBiteBlueProjection}(\omega,u)\bigr)\in C_B\times P_B
\]
is present in \(\noderef{TwoBiteBlueGraph}(\omega)\) and is counted by
\(E^\rightarrow(C)\).

It remains to bound how many incidences can charge to one fixed ordered
candidate pair.  Fix an ordered red pair \((r,s)\in C_R\times P_R\).  A red
incidence charges to \((r,s)\) only when its center is exactly
\(\operatorname{inl}(r)\) and its vertex \(u\) satisfies
\(\noderef{TwoBiteRedProjection}(\omega,u)=s\).  Thus all such vertices
\(u\) lie in the red fiber \(\noderef{RedFiber}(\omega,s)\).  The red
fiber-size part of \(\noderef{FiberAndDegreeEvent}\) says
\[
\noderef{WithinRelativeError}\!\left(
  \varepsilon_2,(\log n)^2,|\noderef{RedFiber}(\omega,s)|\right).
\]
By the definition of \(\noderef{WithinRelativeError}\), the last display
implies
\[
|\noderef{RedFiber}(\omega,s)|
\le (1+\varepsilon_2)(\log n)^2
\le L,
\]
where the final inequality is the assumed charge-factor inequality.  Hence
the ordered red pair \((r,s)\) receives at most \(L\) red incidences.  The same
argument, using the blue fiber-size part of \(\noderef{FiberAndDegreeEvent}\)
and \(\noderef{BlueFiber}\), shows that every ordered blue candidate pair
receives at most \(L\) blue incidences.

Summing the preceding multiplicity bound over all ordered present red and
blue candidate pairs counted by \(E^\rightarrow(C)\) proves the charge
inequality (6).  This is the point at which the oriented count is useful: the
two possible directions \((r,s)\) and \((s,r)\) are different ordered pairs,
and each fixed ordered pair uses only one projection fiber.

Now apply \(\noderef{MediumClosedPairsDeterministicReduction}\) with
\[
M=M_I,\qquad A=A_0,\qquad
t_2=\noderef{TwoBiteLargeCutoff}(\varepsilon,n),\qquad
\ell=\noderef{TwoBiteSmallCutoff}(\varepsilon,n),\qquad
L=3(\log n)^2,
\]
with \(w(x)=|\noderef{TwoBiteX}(\omega,I,x)|\) and with
\(E(C)=E^\rightarrow(C)\).  Its hypotheses are exactly: \(A_0\ge0\), \(L>0\),
the large total mass (5), the upper and lower medium bounds (2), and the
charge inequality (6).  The conclusion gives a nonempty \(B\subseteq M_I\)
such that
\[
A_0<\sum_{x\in B}w(x)\le A_0+t_2,\qquad
|B|\ell\le A_0+t_2,\qquad E^\rightarrow(B)>A_0/L. \tag{7}
\]

It remains to convert the ordered count in (7) into a count over genuine graph
coordinates, because the later product-law tail is applied to independent
base-edge coordinates.  In one fixed color, normalize a present ordered pair
\((x,y)\) to \((x,y)\) if \(x<y\), and to \((y,x)\) otherwise.  Since the pair
is present in a simple graph, \(x\ne y\), so the normalized pair is a nonloop
base-edge coordinate.  A fixed normalized coordinate \(\{x,y\}\) has at most
two ordered preimages, namely \((x,y)\) and \((y,x)\).  Therefore the number
of ordered present red candidates is at most twice the number of distinct
normalized red coordinates, and the same statement holds in blue.  Adding the
two colors gives
\[
E^\rightarrow(B)\le 2E^{\mathrm{coord}}_{I,B}.
\]
Together with (7) and \(L>0\), this implies
\[
E^{\mathrm{coord}}_{I,B}>A_0/(2L)=A_0/(6(\log n)^2).
\]
Substituting the definitions of \(w,t_2,\ell,L\), and unpacking
\(B\subseteq M_I\), gives the claimed \(I\), the claimed nonempty medium
family \(B\), the two size controls, and the normalized graph-coordinate lower
bound.
\end{proof}
